\subsection{Domain and Range from a Graph} 
If all we know about a relation is its graph, we can still find the domain and range.
\begin{example}
The relation $R$ is given by the graph below.
\begin{icpic}[x=0.5\linespace,y=0.5\linespace]
\setgraphsquare{5}
\GraphSmall
\draw[graphend] (-3,-3) parabola bend (-1,1) (0,0);
\draw[graphend] (-1,2) arc (180:0:2);
\draw[dotempty] (-1,2) circle (0.2);
\draw[dotempty] (-3,-3) circle (0.2);
\draw[dotempty] (3,2) circle (0.2);
\draw[dotfull] (0,0) circle (0.2);
\draw[blue] (5,5) node[anchor=north east]{$R$};
\end{icpic}
\begin{itemize}
\item To find the domain of $R$:
\begin{itemize}
\item Imagine we ``flattened'' the graph into the \makeblanksmall-axis.
\item $\dom R=\makeblank$
\end{itemize}
\item To find the range of $R$:
\begin{itemize}
\item Imagine we ``flattened'' the graph into the \makeblanksmall-axis.
\item $\rng R=\makeblank$
\end{itemize}
\end{itemize}
\end{example}
\begin{note}
If the graph continues off the edge of the grid, we assume it extends indefinitely.
\end{note}
\begin{example}
The relation $S$ is graphed below.
\begin{icpic}[x=0.5\linespace,y=0.5\linespace]
\setgraphsquare{5}
\GraphSmall
\draw[graph,rotate=90] (3,-5) parabola bend (-1,2) (-5,-5);
\draw[blue] (5,5) node[anchor=north east]{$S$};
\end{icpic}
\begin{itemize}
\item To find the domain of $S$:
\begin{itemize}
\item Imagine we ``flattened'' the graph into the \makeblanksmall-axis.
\item The domain extends indefinitely \makeblank
\item $\dom R=\makeblank$
\end{itemize}
\item To find the range of $S$:
\begin{itemize}
\item Imagine we ``flattened'' the graph into the \makeblanksmall-axis.
\item The range extends indefinitely \makeblank
\item $\rng R=\makeblank$
\end{itemize}
\end{itemize}
\end{example}